% discussion.tex
We have shown that pattern multiplicity in the modern Hopfield energy provides a principled,
training-free mechanism for conditioning protein sequence generation on functional subsets. The
theoretical framework is exact: multiplicity weights tilt the Boltzmann distribution through a
simple logit bias, and the resulting Langevin dynamics sample from a known distribution with no
approximation. A single scalar $\rho$ continuously interpolates between unconditioned generation
and hard subset curation, and the practitioner can compute the required $\rho$ for any target
effective designated fraction directly from the family composition (Eq.~\ref{eq:rho-inverse}). On
the Kunitz domain, curated memory achieved 100\% phenotype transfer with as few as three
designated input sequences; multiplicity-weighted generation on the full family increased
phenotype fidelity from the natural proportion (32\%) to 63\% at $\rho = 500$; and the
cross-family validation on
Kunitz, SH3, and WW domains showed a quantitative relationship between PCA-space geometry and
conditioning effectiveness.

The central finding is the calibration gap and its decomposition. The gap is not a
failure of the multiplicity-weighted energy---the attention weights track $f_{\mathrm{eff}}$ with
$< 0.3\%$ deviation across all families and $\rho$ values tested---but a well-characterized
property of the PCA encoding that compresses $20L$-dimensional sequence space into $d$ principal
components. When the designated and background subsets overlap in PCA space (low $S$), the
energy-level tilt does not propagate to the decoded residue identity; when they separate
(high $S$), it does. The cross-family validation produced a linear relationship between $S$ and
$\Delta$ with $R^{2} \approx 1.0$, suggesting that the separation index is a reliable \emph{a
priori} predictor of conditioning effectiveness. This yields a practitioner's decision rule:
compute $S$ (a cheap operation requiring only PCA and pairwise cosine similarities, no
generation); if $S > 0.3$, multiplicity weighting preserves fold-level constraints from the full
family while achieving high phenotype transfer (SH3 domains); if $S < 0.15$, hard
curation is necessary because the PCA bottleneck attenuates the signal too severely (WW domains);
the intermediate regime ($0.15 < S < 0.3$, Kunitz domains) admits a tradeoff controlled by
$(\rho,\, \beta/\beta^{*})$.

Hard curation achieves complete phenotype transfer with as few as three designated
input sequences. Many experimental campaigns---phage display, yeast
surface display, functional selections, or even literature curation---produce exactly this kind of
small, functionally characterized set. The practitioner brings external knowledge (``these
sequences bind my target,'' ``these variants are thermostable,'' ``these hits came from my
screen'') and SA converts this small set into a large, diverse candidate library while preserving
the input phenotype. This ``small-set amplifier'' capability complements high-throughput screening:
the screen identifies a few hits, SA expands them computationally, and the expanded library can be
fed back into a second round of selection or directly into functional assays. The method itself is
agnostic to the meaning of the designation---it biases the sampler toward the designated patterns
regardless of whether the underlying phenotype is binding, stability, expression, or any other
property. The scaling study shows that even $K_{\mathrm{des}} = 3$ suffices for phenotype
transfer; diversity improves with more input sequences but saturates near
$K_{\mathrm{des}} = 20$.

The $\omega$-conotoxin experiment extends the small-set amplifier result to a therapeutically
relevant peptide family and illustrates a practically important design choice: whether to seed SA
from the full family or from a curated functional subset. Full-family seeding is appropriate when
the goal is to generate novel family members with broad structural diversity; curated seeding is
necessary when the goal is to preserve a specific pharmacophore. In the $\omega$-conotoxin case,
the O-superfamily is structurally coherent (all members share the C--C--CC--C--C disulfide
framework) but functionally heterogeneous (members block N-type, P/Q-type, or sodium channels
with varying selectivity). Full-family seeding dilutes the Cav2.2-specific Tyr pharmacophore
because the PCA encoding reflects the full structural diversity; curated seeding from the 23
N-type blockers concentrates the PCA subspace on Cav2.2-relevant variation, allowing SA to
regenerate it with 98\% fidelity. This generalizes the separation-index interpretation: even
when $S$ is moderate, hard curation can achieve high pharmacophore transfer by
restricting the memory to the functional subset of interest. The cost---lower sequence diversity
and reduced novelty (0.46 vs.\ 0.60)---reflects the narrower PCA basis, consistent with the
Kunitz scaling study.

The AlphaFold2-multimer docking analysis provides a limited computational check on binding
plausibility. Complex-level confidence scores (pTM $\sim 0.22$, iPTM $\sim 0.10$, interface
pLDDT $\sim 38$--$41$) were modest and overlapping across SA-generated and natural control
sequences, with no pairwise differences reaching significance. These values are consistent with
retained pharmacophore plausibility but do not constitute evidence of binding; larger control
sets and biochemical validation are needed to assess whether pharmacophore inheritance translates
to functional activity.

Several limitations should be acknowledged. The cross-family Pfam regression rests on three points
(with a fourth, $\omega$-conotoxin, following the monotonic trend but lying off the Pfam line);
additional families are needed to test the linearity of the $S$--$\Delta$ relationship and to
assess its generality beyond single-position phenotypic markers. Our functional splits are
defined by single marker positions (P1 in Kunitz, Trp in SH3, a specificity-loop residue in WW);
multi-position phenotypes---where binding specificity depends on coordinated residue identities at
several positions---may exhibit different calibration behavior. We have not experimentally
validated the binding activity of generated sequences; the phenotype fidelity reported here is a
necessary but not sufficient condition for functional binding. For the $\omega$-conotoxin family,
we validated SA-generated sequences against published SAR data from alanine-scanning and
site-directed mutagenesis studies~\cite{kimTyr13Essential1995,satoBasicResidues1993,lewisConotoxinSAR2012}
and found that strong-seeded generation preserved all experimentally identified binding
determinants (Table~\ref{tab:sar-agreement}); for the Kunitz domain, deep mutational scanning
data for the BPTI--trypsin interaction are now
available~\cite{heyneBPTIDMS2021} and could provide a quantitative binding-landscape comparison
in future work. Computational binding affinity predictions (e.g., via AlphaFold-Multimer or
FoldX) could provide an additional intermediate validation tier between sequence-level markers
and experimental assays. Finally, we tested interface-weighted PCA
encoding as a strategy to close the calibration gap, but found it neutral at low weights and
destructive at high weights; the PCA bottleneck reflects the geometric entanglement of functional
and structural variation in the family, not the absence of interface-position components.

Looking forward, the multiplicity framework opens several directions. The binary split
(designated/background) generalizes naturally to multi-class conditioning by assigning different
multiplicities to different functional categories. The PCA bottleneck that limits conditioning
effectiveness might be reduced by alternative sequence encodings---learned representations from
protein language models, for instance, may better separate functional subsets than linear PCA.
The $\beta$ lever suggests that adaptive temperature schedules, increasing $\beta$ during sampling
to sharpen the multiplicity bias, could improve phenotype transfer without permanently sacrificing
diversity. These extensions remain within the analytic, training-free framework that makes
stochastic attention practical for the scarce-data regime.
